% ===== cheat-sheet.tex: AI in Industry, HS26. The semester on one A3 page, drawn as a mind map =====
% Carlos, 2026-11-11: "produce a cheat sheet that, if possible, acts as a mind map of all the concepts ... that
% are being covered in the semester, and that acts as an orientation guide for the participants on when to use
% each technology and what is the main use and application in industry."
%
% The concepts are the lecture rows of the master spreadsheet's Syllabus tab and of the course site's
% src/data/weekends.ts (both read 2026-11-11), in the words of the course's decks. Every named company, product
% and year on the page has its source in SOURCES.md, and no example is older than November 2021.
%
% The look is the course's Codex plates (../style/CLAUDE.md): the colours, the type and the hatching below are
% copied from ../style/ainit-preamble.tex. The page is a plain article in A3 landscape, because the preamble is
% a beamer theme for 160 x 90 mm frames.
%
% Build and check from this folder:
%   ../style/tools/build.sh cheat-sheet.tex
%   ../style/tools/check.sh cheat-sheet.pdf
\documentclass{article}
\usepackage[paperwidth=420mm, paperheight=297mm, margin=0mm]{geometry}
\usepackage{fontspec}
\usepackage{xcolor}
\usepackage{tikz}
\usetikzlibrary{calc, patterns.meta, backgrounds}
\pagestyle{empty}
\setlength{\parindent}{0pt}

% ----- type: copied from ../style/ainit-preamble.tex -----
\setmainfont{EBGaramond}[
  Extension      = .otf,
  UprightFont    = *-Medium,
  ItalicFont     = *-MediumItalic,
  BoldFont       = *-SemiBold,
  BoldItalicFont = *-SemiBoldItalic,
  Ligatures      = TeX,
  Numbers        = Lining,
  Scale          = 1.08]
\providecommand{\geezfontpath}{/System/Library/Fonts/Supplemental/}
\newfontfamily\geezlightfont[Path=\geezfontpath, RawFeature={instance=Light}]{KefaIII.ttf}

% ----- colours: copied from ../style/ainit-preamble.tex -----
\definecolor{paper}{HTML}{F9F4E3}
\definecolor{ink}{HTML}{241F1A}
\definecolor{muted}{HTML}{5E5448}
\definecolor{brand}{HTML}{58409A}
\definecolor{accent2}{HTML}{8F5A06}
\colorlet{ruling}{ink!40!paper}
\colorlet{ruledark}{ink!62!paper}
\colorlet{nib}{ink!88!paper}
\pagecolor{paper}
\color{ink}

% ----- the drawing kit of this page -----
\tikzset{
  cx page/.style={remember picture, overlay, shift={(current page.north west)}, x=1mm, y=-1mm},
  hatch fine/.style={pattern={Lines[angle=45, distance=0.62mm, line width=0.36pt]}, pattern color=brand},
  arm/.style={draw=brand, line width=1.1pt, line join=miter},
  tree/.style={draw=nib, line width=0.5pt, line join=miter},
  card/.style={draw=nib, line width=0.5pt, line join=miter, fill=paper, anchor=north west,
    inner xsep=1.8mm, inner ysep=1.5mm, align=flush left, text=ink,
    font=\fontsize{8.4}{10.1}\selectfont},
  group/.style={anchor=north west, inner xsep=0mm, inner ysep=0.5mm, text=ink, text height=2.75mm, text depth=0.85mm,
    font=\fontsize{10.5}{12}\selectfont\bfseries\itshape},
}
\newcommand{\lab}[1]{{\itshape\color{muted}#1}\enspace}

% ----- the tree of one column -----
% \cxcolumn{spine x}{group x}{card x}{card text width}{start y}{spine top y}: the spine runs down the column,
% a group sits on it with a hatched square, and its cards hang off it on short level branches.
\newcommand{\cxgroupgap}{3.9}
\newcommand{\cxcardgap}{1.9}
\newcommand{\cxcolumn}[6]{%
  \coordinate (sx) at (#1,0);
  \coordinate (gx) at (#2,0);
  \coordinate (cx) at (#3,0);
  \gdef\cxtw{#4}%
  \coordinate (cxstart) at (#3,#5);
  \gdef\cxprev{cxstart}%
  \gdef\cxtop{#6}}
% \cxgroup{node}{label}
\newcommand{\cxgroup}[2]{%
  \path (gx |- \cxprev.south) ++(0,\cxgroupgap) node[group] (#1) {#2};
  \gdef\cxprev{#1}}
% \cxcard{node}{name}{what it is}{use it when}{in industry}; an empty field is left out
\newcommand{\cxcard}[5]{%
  \path (cx |- \cxprev.south) ++(0,\cxcardgap) node[card, text width=\cxtw mm] (#1)
    {{\color{brand}\bfseries #2}\enspace #3%
     \if\relax\detokenize{#4}\relax\else\par\vspace{0.5mm}\lab{Use it when}#4\fi
     \if\relax\detokenize{#5}\relax\else\par\vspace{0.5mm}\lab{In industry}#5\fi};
  \begin{scope}[on background layer]\draw[tree] (sx |- #1.west) -- (#1.west);\end{scope}
  \gdef\cxprev{#1}\gdef\cxlast{#1}}
% \cxspine: the spine of the column, from its top to the branch of the last card
\newcommand{\cxspine}{%
  \begin{scope}[on background layer]\draw[tree] (sx |- 0,\cxtop) -- (sx |- \cxlast.west);\end{scope}}
% \cxmark{group node}: the hatched square where a group sits on the spine
\newcommand{\cxmark}[1]{%
  \path (sx |- #1.west) coordinate (cxm);
  \fill[paper] ($(cxm)+(-1.45,-1.45)$) rectangle ($(cxm)+(1.45,1.45)$);
  \fill[hatch fine] ($(cxm)+(-1.45,-1.45)$) rectangle ($(cxm)+(1.45,1.45)$);
  \draw[brand, line width=0.55pt, line join=miter] ($(cxm)+(-1.45,-1.45)$) rectangle ($(cxm)+(1.45,1.45)$);}

% ----- the head of one weekend -----
% \cxhead{spine x}{text x}{numeral}{Weekend N, dates}{title}{what it is for}{project}
\newcommand{\cxhead}[7]{%
  \draw[brand, line width=1.5pt, line join=miter] (#1-6,43) rectangle (#1+6,55);
  \draw[brand, line width=0.35pt] (#1-4.6,44.4) rectangle (#1+4.6,53.6);
  \node[text=brand, inner sep=0, font=\fontsize{23}{23}\selectfont] at (#1,49.2) {{\geezlightfont #3}};
  \node[anchor=base west, inner sep=0, text=muted, font=\fontsize{9.5}{11}\selectfont\itshape] at (#2,46.4) {#4};
  \node[anchor=base west, inner sep=0, text=ink, font=\fontsize{14.5}{16}\selectfont\bfseries] at (#2,53.8) {#5};
  \node[anchor=base west, inner sep=0, text=ink, font=\fontsize{10}{12}\selectfont\itshape] at (#2,59.9) {#6};
  \node[anchor=base west, inner sep=0, text=ink, font=\fontsize{9.5}{11}\selectfont] at (#2,65.1) {\lab{Project}#7};}

\begin{document}
\null
\begin{tikzpicture}[cx page]

% ===== the leaf: a fine double ruled block, pricking dots in the outer margin =====
\begin{scope}[on background layer]
  \draw[ruling, line width=0.9pt] (10,10) rectangle (410,287);
  \draw[ruling, line width=0.4pt] (11.6,11.6) rectangle (408.4,285.4);
  \foreach \i in {0,...,24} {\fill[ruling] (5.2,{10+11.5417*\i}) circle (0.5);}
\end{scope}

% ===== the root and its three arms =====
\node[anchor=base, inner sep=0, text=ink, font=\fontsize{30}{32}\selectfont\bfseries] at (210,27.2) {AI in Industry};
\draw[ruledark, line width=0.9pt] (183,31.2) -- (237,31.2);
\draw[ruling, line width=0.35pt] (183,32.5) -- (237,32.5);
\draw[arm] (210,32.5) -- (210,37.6);
\draw[arm] (20,43) -- (20,37.6) -- (327.2,37.6) -- (327.2,43);
\draw[arm] (173.6,37.6) -- (173.6,43);

% ===== the three heads =====
\cxhead{20}{29.6}{፩}{Weekend 1, 16 and 17 October}{Large language models}
  {Tasks in words: answer, write, summarize, act}{LLM marketing agent}
\cxhead{173.6}{183.2}{፪}{Weekend 2, 13 and 14 November}{Representation learning}
  {Matching: which product for which customer}{Recommender for online retail}
\cxhead{327.2}{336.8}{፫}{Weekend 3, 11 and 12 December}{Reinforcement learning}
  {Decisions in sequence, rewards later}{RL vehicle routing}
% the second column of weekends 1 and 2 branches off under the head
\draw[tree] (20,69) -- (93.8,69);
\draw[tree] (173.6,69) -- (247.4,69);

% ===== weekend 1, first column =====
\cxcolumn{20}{25}{26}{57.2}{69}{55}
\cxgroup{g11}{Foundations}
\cxcard{lm}{Language model}{A transformer gives each possible next word a chance; one is drawn, then the next.}
  {the input and the output are text.}{chat assistants, drafting, summaries, translation.}
\cxcard{emb}{Embedding}{A text as a row of numbers: similar meanings, similar rows.}
  {comparing or searching texts by meaning.}{search by meaning; the search step of RAG.}
\cxgroup{g12}{Prompting and retrieval}
\cxcard{pe}{Prompt engineering}{A brief with role, task, examples and format; chain or tree of thought for reasoning in steps.}
  {starting any task: it trains nothing.}{customer replies, marketing texts, sorting messages.}
\cxcard{rag}{RAG}{Retrieval-augmented generation: search your own documents, then answer from the pages found.}
  {answers must rest on your documents and cite them.}{DoorDash's support bot answers couriers from its help articles (2024).}
\cxgroup{g13}{Prompt optimization}
\cxcard{opro}{OPRO}{Optimization by prompting: an LLM reads earlier prompts with their scores and writes better ones.}
  {you have a score and examples to compute it on.}{tuning a prompt to a metric, with no retraining.}
\cxcard{gepa}{GEPA}{An LLM reads why a prompt failed and rewrites it; every prompt that is best on some example is kept.}
  {examples are few and each run is costly.}{Databricks: an open model with its prompts beat costlier models at extracting data from documents (2025).}
\cxspine
\cxmark{g11}\cxmark{g12}\cxmark{g13}

% ===== weekend 1, second column =====
\cxcolumn{93.8}{98.8}{99.8}{57.2}{69}{69}
\cxgroup{g14}{Training}
\cxcard{sft}{Supervised fine-tuning}{Train the model further on examples of good answers written by people.}
  {a prompt cannot give the manner or format you need.}{chat assistants: InstructGPT (OpenAI, 2022).}
\cxcard{rlhf}{Human feedback (RLHF)}{People rank answers, a scorer learns their taste, the model is nudged toward high scores.}
  {a good answer is a matter of taste.}{ChatGPT (OpenAI, 2022).}
\cxcard{rlvr}{Verifiable rewards (RLVR)}{A program checks each answer; right answers become more likely.}
  {a program can check the answers: maths, code.}{reasoning models such as DeepSeek-R1 (2025).}
\cxgroup{g15}{Agents}
\cxcard{harness}{Harness}{Software around the model: the model picks a tool, the harness runs it and returns the result, in a loop.}
  {the task needs actions in files, databases or the web.}{coding agents, deep research, support bots that act.}
\cxcard{plan}{Planning strategies}{Try possible steps in a search tree before acting: Monte Carlo tree search, LATS; UCB picks the branch.}
  {a long task punishes an early wrong step.}{AlphaDev found faster sorting code by tree search, now in a standard C++ library (DeepMind, 2023).}
\cxgroup{g16}{Production}
\cxcard{mlops}{MLOps}{Machine learning operations: one pipeline from data to training, evaluation, deployment and monitoring.}
  {a model goes live and must stay reproducible.}{Uber's Michelangelo, one platform for all its teams (2024).}
\cxspine
\cxmark{g14}\cxmark{g15}\cxmark{g16}

% ===== weekend 2, first column =====
\cxcolumn{173.6}{178.6}{179.6}{57.2}{69}{55}
\cxgroup{g21}{Similarity}
\cxcard{ip}{Inner product}{Multiply the matching numbers of two rows and add: one score. Cosine similarity: the same on rows scaled to length one.}
  {scoring a customer against a product.}{the score inside search and recommenders.}
\cxcard{nn}{Nearest neighbors}{Predict from the most similar customers or products.}
  {you need a fast shortlist, or ``customers like you''.}{the first shortlist of search at Etsy (2024) and Airbnb (2025).}
\cxgroup{g22}{Learned profiles}
\cxcard{mp}{Matrix product}{Every customer's row times every product's column: all the scores at once.}
  {scoring every customer against every product.}{what GPUs compute fastest; inside recommenders and LLMs.}
\cxcard{mf}{Matrix factorization}{A few learned numbers per customer and per product; their inner products fill in the missing ratings.}
  {there are many customers and products, few ratings each.}{audiobook recommendations at Spotify, with the numbers learned by networks (2024).}
\cxcard{gd}{Gradient descent}{Change every number a small step that lowers the error; repeat.}
  {a model has many numbers to learn.}{trains recommenders, neural networks and LLMs.}
\cxspine
\cxmark{g21}\cxmark{g22}

% ===== weekend 2, second column =====
\cxcolumn{247.4}{252.4}{253.4}{57.2}{69}{69}
\cxgroup{g23}{Bandits}
\cxcard{mab}{Multi-armed bandit}{Show one option, learn from the click, choose again: explore or exploit.}
  {choices repeat, feedback is fast and the best option is unknown.}{which offer, page or content to show.}
\cxcard{rnd}{Random}{Pick an option at random each time.}
  {you need a baseline, or a fair test.}{the random split of visitors in an A/B test.}
\cxcard{ete}{Explore then exploit}{Try every option for a while, then keep the best.}
  {the best option does not change.}{an A/B test, then the winner for everyone.}
\cxcard{eg}{Epsilon-greedy}{The best option so far; with chance epsilon, a random one.}
  {you want a simple rule that never stops exploring.}{the mix of music and podcasts on Spotify's home page (2025).}
\cxcard{ucb}{UCB}{Upper confidence bound: the highest estimate plus a bonus for uncertainty.}
  {every try is costly.}{Amazon tested it on the application pages of Amazon Payments (2026).}
\cxgroup{g24}{Scale}
\cxcard{pmf}{Parallelized matrix factorization}{Split customers, products and ratings over many computers that work at once.}
  {the data or the model do not fit one computer.}{LinkedIn trains its ranking models on many GPUs at once (2024).}
\cxcard{hpc}{Computing clusters}{Thousands of GPUs joined by fast links; a run costs GPU hours, energy and money.}
  {training or serving large models.}{Apertus, the Swiss open LLM, trained on Alps at CSCS (2025).}
\cxspine
\cxmark{g23}\cxmark{g24}

% ===== weekend 3 =====
\cxcolumn{327.2}{332.2}{333.2}{65.2}{69}{55}
\cxgroup{g31}{Problem}
\cxcard{mdp}{Markov decision process}{States, actions, rewards and the chances of the next state; a policy says what to do in each state.}
  {each decision changes the next situation.}{maintenance, pricing, stock, delivery routes, robots.}
\cxgroup{g32}{Exploration}
\cxcard{eg3}{Epsilon-greedy}{The bandit rule of weekend 2, used in every state.}
  {you learn while you act.}{}
\cxgroup{g33}{Values}
\cxcard{bell}{Bellman equation}{The value of an action: its reward now plus the value of the best next action.}
  {the rules and their chances are known.}{planning maintenance and stock with known odds.}
\cxcard{rm}{Robbins-Monro}{Learn an average from noisy samples: move the guess a small step toward each one.}
  {every observation is noisy.}{the update inside SARSA, Q-learning and stochastic gradient descent.}
\cxgroup{g34}{Learning from experience}
\cxcard{sarsa}{SARSA and Q-learning}{Learn the value of each action from your own steps, with no model of the world.}
  {there is no model, only experience or a simulator.}{Lyft matches riders and drivers by future earnings learned from live trips (2024).}
\cxcard{rf}{REINFORCE}{Make the actions of well-paid episodes more likely: learn the policy itself.}
  {actions are many, or the policy is the goal.}{the family that trains chat and reasoning models of weekend 1.}
\cxspine
\cxmark{g31}\cxmark{g32}\cxmark{g33}\cxmark{g34}

\end{tikzpicture}
\end{document}
